\documentclass{article}
\usepackage{amsmath}
\usepackage{graphicx}
\usepackage{float} % Required for [H]
\usepackage{hyperref}
\usepackage[en-GB]{datetime2}
\usepackage{tabularx}
\DTMlangsetup[en-GB]{ord=omit}
\graphicspath{{./images}}
\hypersetup{
    colorlinks=true,
    linkcolor=blue,
    urlcolor=blue
}


\title{Epidemiology - Week 1 - Introduction to Epidemiology}
\author{Fernaldy}
\date{\today}


\begin{document}
\maketitle

\section*{What is Epidemiology?}
It is the study of the \textbf{frequency and distribution of health-related states or events in specific populations},
including the study of determinants influencing such states, and the application of this study to control health problems.
\\~\\
Basically, studying the frequency or distribution of health-related outcomes within specific populations
alongside its causation and correlation with the goal of translating the insight into public health actions.

\subsection*{Role of Epidemiology}
\begin{itemize}
    \item \textbf{DESCRIBE} - the patterns of health outcome and disease within populations
    \item \textbf{INTERPRET} - the differences
    \item \textbf{APPLY} - the results to public health practice
    \item \textbf{EVALUATE} - the effect of the application / intervention
\end{itemize}

\subsection*{Cases and Populations}
Quantification needs a clear case definition and a defined study population.
\begin{itemize}
    \item \textbf{CASE}: A set of standard criteria to decide whether a person has a specific disease, syndrome, or health condition
    \item \textbf{STUDY POPULATION}: The population who could become or have become cases, AKA \textbf{population at risk}
\end{itemize}

\subsection*{Exposure and Outcome}
\begin{itemize}
    \item \textbf{EXPOSURE}: Risk factor or determinant -> Any factor that may influence the outcome
    \item \textbf{OUTCOME}: Disease, event, or health related state of interest, or risk factor / precursor to a disease
\end{itemize}

\subsection*{Types of Epidemiological Investigation / Study Design}
\begin{figure}[htbp]
    \centering
    \includegraphics[width=\textwidth]{w1-1.png}
    \caption{Types of epidemiological investigation}
    \label{fig:w1-1}
\end{figure}


\subsubsection*{Observational Epidemiology}
This type of epidemiology is based on examining the distribution and determinants of an outcome \textbf{without any
attempt to influence it}, and is split into two further subtypes.
\begin{itemize}
    \item \textbf{DESCRIPTIVE}: Designed to describe the distribution of an exposure or outcome without attempting to explain (or influence) the distribution -$>$ \textit{Who, what, were, when?}
    \item \textbf{ANALYTICAL}: Designed to examine \textbf{association} with the aim of identifying causes for an outcome -$>$ \textit{Why?}
\end{itemize}

\pagebreak
\section*{Outbreaks and Epidemics}
Some definitions of disease occurrence

\begin{figure}[H]
    \centering
    \includegraphics[width=0.6\textwidth]{w1-2.png}
    \caption{Disease occurrence types}
    \label{fig:w1-2}
\end{figure}

\begin{itemize}
    \item \textbf{ENDEMIC}: Common at baseline level in a given area / population
    \item \textbf{OUTBREAK}:
        \begin{itemize}
            \item Excess of cases of a disease compared with \textit{expected} number
            \item Limited to a local increase in disease (often infectious) incidence; e.g. Town, village, institution
        \end{itemize}
    \item \textbf{EPIDEMIC}:
        \begin{itemize}
            \item Excess of cases of a disease compared with \textit{expected} number
            \item Generalised, cases distributed across wider area / time period
        \end{itemize}
    \item \textbf{PANDEMIC}: Affects large numbers of people and crosses international boundaries
\end{itemize}

\subsection*{Examples of Types of Epidemics}

\begin{table}[H]
\centering
\begin{tabularx}{\textwidth}{|l|X|X|}
\hline
\textbf{Type}          & \textbf{Description}                                             & \textbf{Example}              \\ \hline
\textbf{Common-source} & Individuals are exposed to the same source of infection / hazard & Food contamination at wedding \\ \hline
\textbf{Propagated}    & Disease spreads from person-to-person                            & Measles outbreak at school    \\ \hline
\textbf{Mixed}         & A combination of common-source followed with propagation          & Cholera outbreak that started from a contaminated water source followed with household transmission \\ \hline
\end{tabularx}
\caption{Example of types of epidemics}
\label{tab:example-epi}
\end{table}

\subsection*{Clusters}
An aggregation of some relatively uncommon disease or event which can include infectious and non-communicable
conditions. Often unknown / multiple aetiological factors with low exposure level over number of years.


\pagebreak
\section*{Pandemics}
A pandemic is an epidemic that \textit{spreads across multiple countries or continents} and \textit{affects large numbers of people}, \textit{involving a novel pathogen or population susceptibility}. Examples,
\begin{itemize}
    \item 1918-1919 Influenza pandemic
    \item Covid-19
\end{itemize}

\subsection*{Strategy against Pandemics}
\begin{figure}[H]
    \centering
    \includegraphics[width=1.1\textwidth]{w1-3.png}
    \caption{Pandemic strategy choices}
    \label{fig:w1-3}
\end{figure}

\subsubsection*{Elimination vs Mitigation}
Eliminate literally just means trying to eliminate / stop the community transmission of the disease, while mitigation
aims to slow the transmission down to a manageable level so healthcare systems don't get overwhelmed.


\end{document}
